\section{Event-time derivatives and hybrid sensitivity}\label{sec:hybrid}
\subsection{One event at a fixed comparison time}
Before an event let $\dot z=f_a^-(t,z)$; let $g(a,t,z)=0$ trigger reset $z^+=\mathcal R(a,t,z^-)$; afterward let $\dot z=f_a^+(t,z)$. At the nominal event $\tau$, put
\[
 d_g=g_t+D_zg\,f^-\ne0.
\]
The incoming sensitivity $s^-$ is evaluated at the fixed nominal time, whereas the moving event state has total derivative $s^-+f^-\tau'$.

\begin{proposition}[Event-time and parameter saltation]\label{prop:saltation}
\begin{align}
 \tau'&=-\frac{g_a+D_zg\,s^-}{d_g},\\
 s^+&=D_z\mathcal R\,s^-+\mathcal R_a
 +(D_z\mathcal R f^-+\mathcal R_t-f^+)\tau'.
 \label{eq:saltation}
\end{align}
Equivalently, the homogeneous sensitivity multiplier is
\[
 D_z\mathcal R+
 \frac{(f^+-D_z\mathcal R f^--\mathcal R_t)D_zg}{d_g},
\]
and the inhomogeneous parameter term is
$\mathcal R_a+(f^+-D_z\mathcal R f^--\mathcal R_t)g_a/d_g$.
\end{proposition}
\begin{proof}
Differentiate $g(a,\tau(a),z_a^-(\tau(a)))=0$ to get
$g_a+D_zg\,s^-+d_g\tau'=0$. The total derivative of the reset event point is
$D_z\mathcal R(s^-+f^-\tau')+\mathcal R_a+\mathcal R_t\tau'$.
To compare at the same absolute time, subtract the post-event displacement $f^+\tau'$. This yields the second formula. Substitute the first to obtain the multiplier and parameter term. Omitting the comparison-time correction would be an incorrect derivative of the hybrid trajectory.
\end{proof}

This derivation is consistent with standard saltation analysis \cite{KongEtAl}; the proof is included to fix all signs and parameter terms.

\subsection{Second, third, and arbitrary finite order}
Write $F(a,t)=g(a,t,z_a^-(t))$. Since $F(a,\tau(a))=0$,
\begin{align}
 \tau''&=-\frac{F_{aa}+2F_{at}\tau'+F_{tt}(\tau')^2}{F_t},\\
 \tau'''&=-\frac{1}{F_t}\bigl[F_{aaa}+3F_{aat}\tau'+3F_{att}(\tau')^2
 +F_{ttt}(\tau')^3\notag\\
 &\hspace{40mm}+3(F_{at}+F_{tt}\tau')\tau''\bigr].
 \label{eq:tau-high}
\end{align}
Each formula follows by differentiating the preceding identity and collecting the highest derivative of $\tau$. At order $r$, the only new highest term is $F_t\tau^{(r)}$; all other terms are known derivatives and lower jets. Appendix~\ref{app:jets} gives a coefficient-extraction algorithm that determines the entire finite polynomial without an unspecified remainder symbol.

\begin{assumption}[A regular finite-event chamber]\label{ass:chamber}
Fix a common compact initial-state set, compact parameter interval, and finite horizon. The relevant trajectories have a fixed event itinerary of at most $M$ events; $|d_g|\ge\kappa>0$ at each event; event times, nonchosen guards, and the initial and terminal times have positive separation margins. Flights, guards, and resets are $C^{r+1}$ on a common compact neighborhood, with bounded derivatives through the required order. These are dynamical assumptions; regularity of an integrated preparation is specified separately below.
\end{assumption}

\begin{theorem}[Strong chamber response with an explicit preparation]\label{thm:chamber}
Under Assumption~\ref{ass:chamber}, event times, one-sided event states, and the terminal state are $C^r$ jointly in the physical parameter and initial state, with uniform bounds. A fixed smooth finite-dimensional report of these quantities has the same regularity.

For preparation-integrated response assume either a fixed finite measure, or
\[
 \mu_a=Z_a^{-1}(X_a)_\#(w_a\lambda),\qquad
 Z_a=\int w_a\,d\lambda\ge z_*>0,
\]
where $\lambda$ is fixed and finite, $w_a\ge0$, and $X_a$ takes values in the common chamber. The transports and weights are jointly measurable and $C^r$ in the parameter for $\lambda$-almost every source point. Require common integrable envelopes for all derivative products through order $r$ of the composed weighted tests. Then the preparation integral is $C^r$, retaining all transport, weight, observable, and normalization derivatives. For a compact $d$-dimensional report, the analogous envelopes for weighted report jets give strong $C^r$ response in $H^{-s}$ for $s>d/2+r$.
\end{theorem}
\begin{proof}
The first flight has the joint flow regularity obtained from its variational ODE, as in Lemma~\ref{lem:flow}. The implicit function theorem with $|d_g|\ge\kappa$ gives the first event time, jointly in parameter and initial state. Equations~\eqref{eq:saltation}--\eqref{eq:tau-high} and the finite-jet recursion supply each derivative. Compose with the reset and repeat at most $M$ times. Separation of guards prevents a hidden change of itinerary. Every division has denominator bounded away from zero, and all other expressions are finite polynomials of bounded derivatives. This proves the dynamical conclusion independently of any preparation measure.

Let $Y_a(z)$ be the report from $X_a(z)$. The unnormalized preparation integral is
$J_a=\int w_a(z)F_a(Y_a(z))\lambda(dz)$, with derivative
\[
 J_a'=\int\{\dot w_a F_a(Y_a)+w_a\dot F_a(Y_a)
                   +w_a DF_a(Y_a)\dot Y_a\}\,d\lambda.
\]
Here $\dot Y_a$ includes the initial transport derivative $\dot X_a$. Apply the finite Leibniz and chain rules at higher orders. The assumed envelopes justify dominated differentiation and continuity at every order. Finally differentiate $J_a/Z_a$ with the quotient recursion. A fixed preparation is obtained by taking parameter-independent weights and transport.

For the law, write the unnormalized report distribution as $\int w_a\delta_{Y_a}\,d\lambda$. Lemma~\ref{lem:dirac}, the chain rule and the integrable jet envelopes justify Bochner differentiation. Normalize by $Z_a$. This proves the strong distributional conclusion without treating fixed support as differentiability of the preparation.
\end{proof}

For identity dynamics and the fixed support $\{0,1\}$, the cusp weights
$\mu_a=(1/2+|a|)\delta_0+(1/2-|a|)\delta_1$ are not allowed by the preparation hypothesis. Their nonsmooth expectation does not contradict the trajectory-only conclusion. This resolves finding M3 while retaining transported and weighted preparations, not merely fixed ones.

The constants can grow rapidly with $M,r,\kappa^{-1}$. For instance, if a flight sensitivity multiplies by at most $A_1$ and one event/reset has multiplier and forcing at most $B_1$, a first-order bound obeys
$M_{j,1}\le B_1A_1M_{j-1,1}+B_1(A_1+1)$.
Higher orders follow the explicit polynomial recursion. No polynomial-in-$M$ or grazing-uniform bound is claimed from this chamber argument.

\subsection{A terminal crossing must be assembled}
When $\tau(a,z)=T$ cuts through the preparation, terminal velocity can jump between unreflected and reflected values. For almost every fixed initial state the event indicator is locally constant in $a$, yet its expectation can have a nonzero derivative. Split the preparation by $\tau<T$ and $\tau>T$ and retain the interface term of Theorem~\ref{thm:shape}. The local Lorentz example below calculates this term exactly. The full-preparation tube construction then avoids uniform transversality altogether by changing integration variables, not by declaring the chamber estimate uniform near grazing.
\par
